% Notes of the Montreal seminar, 1970 (cote 129, typed part): archive cote 129, site page 41 (PDF page 42).
% typescript. Draft compared: none.
% AI-transcribed by Claude Opus 5.5 (claude-opus-5-5, adaptive thinking, effort high), Message Batches API, batch mtl129_full (msgbatch_01AGnDKHwy6ATCM7ehUa6263), 2026-11-13; not yet checked by eye.
\scanpage{129/41}
\mspage{\hand{II.13}}

F et V satisfaisant les relations suivantes :
\[ F\cdot\lambda = \lambda^{\sigma}\cdot F \]
\[ \lambda\cdot V = V\cdot\lambda^{\sigma} \]
\[ F\cdot V = V\cdot F = p \ , \]
pour tout $\lambda \in W$ et où l'on note par $\lambda^{\sigma}$ l'image de $\lambda$ par le morphis-%
me de Frobénius de $W$ .

On appelle \uline{catégorie des modules de Dieudonné} la catégorie
$\underline{\mathrm{Mod}}_{\hand{\mathcal{D}};W\ \hand{\ell.f.}}$ des \hand{$\mathcal{D}$}-modules à gauche tels que le W-module sous-jacent soit
de longueur finie. On voit en particulier qu'un tel module est annulé par
une puissance de $p$ .

Un objet de $\underline{\mathrm{Mod}}_{\hand{\mathcal{D}};W\ \hand{\ell.f.}}$ est donc un W-module $M$ de longueur
finie, muni de deux morphismes de W-modules
\[ F_M : M^{\sigma} \rightarrow M \]
\[ V_M : M \rightarrow M^{\sigma} \quad , \]
en posant
\margin{$\times$ \quad $\otimes$}
\[ M^{\sigma} = M \hand{\otimes}_W (W,\sigma) \ , \]
($\sigma$ étant le morphisme de Frobénius de $W$), $F_M$ et $V_M$ vérifiant
\[ F_M\cdot V_M = p\cdot \mathrm{Id}_M \ , \ V_M F_M = p\cdot \mathrm{Id}_M \quad . \]

\uline{Théorème\unsure{.}4.2.} (Dieudonné)

Il existe une équivalence de catégories
\[ D^* : [p\text{-}\mathrm{Gr}\text{-}\mathrm{fin}(k)]^{0} \hand{\xrightarrow{\ \approx\ }} \underline{\mathrm{Mod}}_{\hand{\mathcal{D}},W\ \hand{\ell.f.}} \]
telle que l'on ait un isomorphisme canonique $[D^*(G)]^{\unsure{0|\sigma}} \hand{\approx} D^*(G^{(p)})$ ,

\margin{41}
